\begin{HVTPage}{Capítulo 4}{Biohidrógeno}{Ruta biológica}
\HVTLead{La fermentación oscura utiliza microorganismos para degradar materia orgánica y producir una mezcla gaseosa que puede contener H2 y CO2, además de metabolitos líquidos.}
\HVTEquation{C_6H_{12}O_6+2H_2O\rightarrow2CH_3COOH+2CO_2+4H_2}{}
\HVTText{La reacción representa la ruta estequiométrica teórica hacia acetato, cuyo límite es 4 mol H2/mol glucosa; no representa el rendimiento garantizado de un residuo real. \HVTCite{6} Los rendimientos experimentales dependen del sustrato, la comunidad microbiana, el pH, la temperatura, la carga orgánica y el tiempo de retención. \HVTCite{7}}
\begin{HVTTable}{Variable}{Símbolo}{Unidad}
\HVTTableRow{Demanda química de oxígeno}{DQO}{g O2/L}
\HVTTableRow{pH}{pH}{--}
\HVTTableRow{Temperatura}{$T_b$}{°C}
\HVTTableRow{Tiempo de retención hidráulica}{TRH}{h}
\HVTTableRow{Rendimiento de H2}{$Y_{H_2,s}$}{mol H2/mol sustrato}
\end{HVTTable}
\end{HVTPage}

\begin{HVTPage}{Capítulo 4}{Biohidrógeno}{Caracterización mínima}
\begin{HVTTask}{Caracterización mínima}\item Origen y variabilidad del residuo. \item Materia seca, sólidos volátiles y DQO. \item Carbohidratos disponibles e inhibidores. \item Producción y composición del biogás. \item Ácidos orgánicos y estabilidad del reactor.\end{HVTTask}
\end{HVTPage}
